% multipliers: Z15, Z27, Z28, Z31, Z30; exhaustive search Z1, Z2, Z26, Z33.
% Open questions for the author:
% - Assumed environments theorem, proposition, lemma, corollary, definition, remark (amsthm),
%   and labels from other sections: thm:z8 (self-conjugacy rule, congruences section).
% - \lean{...} is given names relative to the namespace SignedDifferenceSets, with raw
%   underscores; adjust if the macro wants full names or escaped underscores.
% - solutions.md (Z33) names the lift theorem `SignedDifferenceSets.not_hasSDS_of_lift`; in
%   LiftSearch.lean it lies in namespace LiftSearch, so its full name is
%   SignedDifferenceSets.LiftSearch.not_hasSDS_of_lift (used below).
% - Lemma lem:cosetconds (C4, C7: absolute sum, parity, marked counts) may duplicate a lemma in
%   the counting section (Z16); merge if so.
% - Added refs.bib entry arasu-gordon-zhang (Z28's Origin cites their Algorithm 1).
% - The hand argument for N32 after Theorem thm:z2 is ours, written for this paper; it is not
%   in claims/ and not formalized. Drop it if unwanted.

\section{Multipliers and exhaustive search}\label{sec:multipliers}

When $n$ is prime to $v$, McFarland's multiplier theorem makes some translate of an SDS
constant on the orbits of certain automorphisms $x \mapsto tx$, and the same holds for its
coset sums on a quotient whose order is prime to $n$. The search space then shrinks to one
value per orbit, and a computer can exhaust it. This section gives the multiplier theorem in
the form we use (one translate fixed by all multipliers at once), four rules built on it
(Z15, Z27, Z28, Z31), a bound derived from it (Z30), and the cases settled by plain exhaustive
search (Z1, Z2, Z26, Z33). Section~\ref{sec:searchcheck} explains how Lean checks the searches.

\subsection{Notation}
For a finite abelian group $H$ we work in the group ring $\mathbb{Z}[H]$ and identify a
function $a \colon H \to \mathbb{Z}$ with $A = \sum_h a(h)[h]$. We write $H$ also for
$\sum_{h \in H} [h]$, and for an integer $t$ we put $A^{(t)} = \sum_h a(h)[th]$, so that
$A^{(-1)} = \sum_h a(h)[-h]$. The identity
\begin{equation}\label{eq:autocorr}
  \sum_{h} a(h)\,a(h+g) = n\,[g = 0] + \mu \quad (g \in H)
  \qquad\Longleftrightarrow\qquad A A^{(-1)} = n + \mu H
\end{equation}
says that $a$ has \emph{autocorrelation} $n[g=0] + \mu$. An SDS $D$ on $G$ has it with
$n = k - \lambda$ and $\mu = \lambda$; its coset sums on a quotient $H$ with fibres of size
$c$ have it with $\mu = \lambda c$ (Lemma~\ref{lem:contraction}). Summing~\eqref{eq:autocorr}
over $g$ gives $s^2 = n + \mu |H|$ for $s = \sum_h a(h)$.

\begin{definition}\label{def:multipliers}
For a finite abelian group $H$ and an integer $n \neq 0$, let $M(H, n)$ be the set of integers
$t$ prime to $|H|$ such that for every prime $p \mid n$ there is $f \ge 0$ with $tx = p^f x$
for all $x \in H$ (that is, $t \equiv p^f \pmod{\exp H}$).
\end{definition}

For $t$ prime to $|H|$ the map $x \mapsto tx$ is an automorphism of $H$; its orbits are the
\emph{orbits of $t$}.

\begin{lemma}\label{lem:cosetconds}
Let $D$ be an SDS $(v, k, \lambda)$ on $G$, $\phi \colon G \to H$ surjective, $c = v/|H|$ and
$a$ the coset sums of $D$. For $x \in \mathbb{Z}$ let $m_c(x) = c$ if $x \equiv c \pmod 2$ and
$m_c(x) = c - 1$ otherwise. Then $\sum_h |a(h)| \le k$, $k - \sum_h |a(h)|$ is even, and
$k \le \sum_h m_c(a(h))$.
\end{lemma}

\begin{proof}
Let $m(h)$ be the number of marked elements in the coset of $h$. Then $a(h)$ is a sum of
$m(h)$ signs, so $|a(h)| \le m(h) \le c$ and $m(h) \equiv a(h) \pmod 2$; hence
$m(h) \le m_c(a(h))$. Sum over $h$, using $\sum_h m(h) = k$.
\end{proof}

\noindent Lean: \lean{IsSDS.sum_abs_contract_le}, \lean{IsSDS.sum_abs_contract_parity},
\lean{marked_le_maxMarked} with \lean{IsSDS.sum_marked}.

\begin{definition}\label{def:cosetconds}
Let $H$ be a finite abelian group and $c, k \ge 1$, $\lambda$, $s$ integers. A function
$a \colon H \to \mathbb{Z}$ \emph{meets the coset conditions} for $(c, k, \lambda, s)$ if
$|a(h)| \le c$ for all $h$, $\sum_h a(h) = s$, $\sum_h a(h)^2 = k + \lambda(c-1)$,
$\sum_h a(h)a(h+g) = \lambda c$ for $g \neq 0$, and the three conditions of
Lemma~\ref{lem:cosetconds} hold.
\end{definition}

By Lemmas~\ref{lem:contraction} and~\ref{lem:cosetconds}, the coset sums of an SDS meet the
coset conditions for $(c, k, \lambda, s')$ with $s' = \pm s$; since all conditions but the sum
are even in $a$, either $a$ or $-a$ meets them for $(c, k, \lambda, s)$.

\subsection{The multiplier theorem with a fixed translate}

\begin{theorem}[McFarland]\label{thm:multiplier}
Let $a \colon H \to \mathbb{Z}$ have autocorrelation $n[g=0] + \mu$ with $n \neq 0$ and
$\gcd(n, |H|) = 1$. Then there is $b \in H$ with $a(tx + b) = a(x + b)$ for all $x \in H$ and
all $t \in M(H, n)$.
\end{theorem}

The theorem is McFarland's (Gordon~\cite[Thm.~5.1]{gordon}, Tan~\cite[Result~3.2]{tan}); the
classical statement gives, for each $t$, a translate fixed by $t$. That one translate serves
every $t$ at once is the observation of McFarland and Rice~\cite{mcfarland-rice}. We give
the proof, since Lean follows it.

\begin{proof}
Fix $t \in M(H, n)$ and let $W = A^{(t)}A^{(-1)} - \mu H$. Write $s = \sum_h a(h)$, so
$s^2 = n + \mu|H|$.

\emph{Step 1: $n$ divides every coefficient of $W$.} It suffices to show $p^e \mid W$ for
every prime power $p^e \,\|\, n$. Let $f$ be as in Definition~\ref{def:multipliers}, so
$A^{(t)} = A^{(p^f)}$.

If $n$ is squarefree ($e = 1$), work in $\mathbb{F}_p[H]$. There $A^{p} = A^{(p)}$
(Frobenius, since $a(h)^p \equiv a(h)$), so
$A^{(t)} \equiv A^{p^f}$ and
\[
  W \equiv A^{p^f - 1} A A^{(-1)} - \mu H = A^{p^f-1}(n + \mu H) - \mu H
  \equiv \mu\,(s^{p^f - 1} - 1)\,H \pmod p ,
\]
using $AH = sH$. If $p \mid \mu$ this vanishes; otherwise $p \nmid s^2 = n + \mu|H|$
(as $p \nmid |H|$), and $s^{p^f - 1} \equiv 1$ by Fermat, since $p - 1 \mid p^f - 1$.

For general $e$, let $\kappa$ be a finite field of characteristic $p$ containing a primitive
$|H|$-th root of unity, and $\Lambda = W(\kappa)$ its ring of Witt vectors: a complete discrete
valuation ring with uniformizer $p$, containing a primitive $|H|$-th root of unity (the
Teichmüller lift of one in $\kappa$), with a Frobenius automorphism $\sigma$ that fixes
$\mathbb{Z}$ and maps every root of unity $u$ of order prime to $p$ to $u^p$. Let $\chi$ run
over the characters $H \to \Lambda^\times$. For $\chi \ne 1$ we have $\chi(H) = 0$,
$\chi(A)\chi(A^{(-1)}) = n$ and $\chi(A^{(t)}) = \sum_h a(h)\chi(h)^{p^f} = \sigma^f(\chi(A))$.
A ring automorphism of a discrete valuation ring with $\sigma(p) = p$ preserves the
valuation, so $\chi(W) = \sigma^f(\chi(A))\,\chi(A^{(-1)})$ has the valuation of $n$, and
$p^e \mid \chi(W)$. For $\chi = 1$, $\chi(W) = s^2 - \mu|H| = n$. Fourier inversion,
$|H|\,W(g) = \sum_\chi \chi(W)\chi(-g)$, gives $p^e \mid |H|\,W(g)$ in $\Lambda$; $|H|$ is a
unit, and an integer divisible by $p^e$ in $\Lambda$ is divisible by $p^e$ in $\mathbb{Z}$.

\emph{Step 2: $W = n[g]$ for some $g \in H$.} Since $x \mapsto tx$ is an automorphism,
$A^{(t)}A^{(-t)} = (n + \mu H)^{(t)} = n + \mu H$, and $A^{(\pm t)}H = A^{(\pm 1)} H = sH$.
Expanding,
\[
  W W^{(-1)} = (n + \mu H)^2 - 2\mu s^2 H + \mu^2|H|H
  = n^2 + 2\mu\,(n + \mu|H| - s^2)\,H = n^2 ,
\]
and $\sum_h W(h) = s^2 - \mu|H| = n$. So $U = W/n \in \mathbb{Z}[H]$ has $\sum_h U(h)^2 = 1$
(the coefficient of $[0]$ in $UU^{(-1)} = 1$) and $\sum_h U(h) = 1$: $U = [g]$.

\emph{Step 3: $A^{(t)} = [g]A$.} Indeed $WA = A^{(t)}(n + \mu H) - \mu sH = nA^{(t)}$, and
$WA = n[g]A$. In coordinates, $a(x) = a(tx - g)$ for all $x$.

\emph{Step 4: one translate for all $t$.} As $\gcd(n, |H|) = 1$ and $s^2 = n + \mu|H|$, $s$ is
prime to $|H|$, so there is $b \in H$ with $sb = \sum_h a(h)h$. The translate
$a'(x) = a(x + b)$ has the same autocorrelation and $\sum_x a'(x)x = \sum_y a(y)(y - b) = 0$.
Steps 1--3 for $a'$ give $g$ with $a'(x) = a'(tx - g)$. Then, substituting $y = tx$,
\[
  0 = t\sum_x a'(x)x = \sum_x a'(tx - g)\,tx = \sum_y a'(y - g)\,y = \sum_z a'(z)(z + g) = sg,
\]
so $g = 0$. Since $b$ does not depend on $t$, $a'$ is fixed by every $t \in M(H, n)$.
\end{proof}

\noindent Lean: \lean{exists_translate_fixed_of_autocorr} (squarefree $n$, Step~1 by
Frobenius, \texttt{Multiplier.lean}) and \lean{exists_translate_fixed_of_autocorr'} (every
$n \neq 0$, Step~1 in Witt vectors, \texttt{MultiplierWitt.lean}). For an SDS:
\lean{IsSDS.exists_translate_fixed}, \lean{IsSDS.exists_translate_fixed'}; for its coset sums
on a quotient $H$ with $\gcd(|H|, n) = 1$: \lean{IsSDS.exists_translate_fixed_contract},
\lean{IsSDS.exists_translate_fixed_contract'}.

\begin{remark}
Step~1 for squarefree $n$ is elementary; we did not find this proof in print, but every step is
standard. The proof for general $n$ is the usual character argument with ideals of
$\mathbb{Z}[\zeta]$ replaced by the discrete valuation ring $\Lambda$; we chose Witt vectors
because Mathlib provides them, and found no earlier formalization of a multiplier theorem.
\end{remark}

\subsection{Orbit exhaustion}

\begin{theorem}\label{thm:z15}
Let $G$ be a finite abelian group of order $v$, and $k$, $\lambda$ integers with
$n = k - \lambda > 0$ squarefree and $\gcd(n, v) = 1$. Let $t \in M(G, n)$ and let $s$ be an
integer with $s^2 = k + \lambda(v - 1)$. Suppose there is no $D \colon G \to \{0, \pm 1\}$ with
$D(tx) = D(x)$ for all $x$, $k$ nonzero values, $\sum_x D(x) = s$ and
$\sum_x D(x)D(x+g) = \lambda$ for all $g \neq 0$. Then $G$ has no $(v, k, \lambda)$ SDS.
\end{theorem}

\begin{proof}
Let $D$ be an SDS. It has autocorrelation $n[g=0] + \lambda$, so by
Theorem~\ref{thm:multiplier} a translate $D'$ is fixed by $t$; $D'$ is an SDS
(Lemma~\ref{lem:translate}) with $\sum D' = \pm s$ (Lemma~\ref{lem:square}). Either $D'$ or
$-D'$, also fixed by $t$, is excluded by the hypothesis.
\end{proof}

\noindent In Lean the group is $\mathbb{Z}_{m_1} \times \dots \times \mathbb{Z}_{m_r}$ and the
condition $t \in M(G, n)$ reads $t \equiv p^f \pmod{m_i}$ for all $i$; the hypothesis on $D$
is decided by a kernel computation (Section~\ref{sec:searchcheck}), as are the arithmetic
conditions (squarefreeness, primality, coprimality, by trial division).

\begin{remark}
Lean: \lean{not_hasSDS_of_orbitRow} (\texttt{OrbitSearch.lean}); the two largest rows, N71 and
N363, through \lean{not_hasSDS_of_orbitSlices} (\texttt{OrbitSearchSplit.lean}). It settles
\zrows{15} rows. Origin: McFarland's multiplier theorem, as in~\cite[Thm.~5.1]{gordon}, with
the fixed translate of~\cite{mcfarland-rice}; Gordon's database already uses orbit exhaustion
(``Orbit Exhaust''); our part is the search for these rows and its Lean proof.
\end{remark}

\begin{theorem}\label{thm:z27}
Theorem~\ref{thm:z15} holds without the hypothesis that $n$ is squarefree.
\end{theorem}

\begin{proof}
The same, with Theorem~\ref{thm:multiplier} for general $n$.
\end{proof}

\begin{remark}
Lean: \lean{not_hasSDS_of_orbitRow'} (\texttt{OrbitSearchWitt.lean}),
\lean{not_hasSDS_of_orbitSlices'}. It settles \zrows{27} rows, all with $n$ not squarefree.
Origin: McFarland's theorem~\cite[Thm.~5.1]{gordon}, \cite[Result~3.2]{tan}, and orbit
exhaustion as in Gordon's database; the Lean proof through Witt vectors is ours.
\end{remark}

\subsection{Multiplier-fixed coset sums}

When the orbits of $t$ on $G$ are too many to search, one searches the coset sums on a
quotient instead.

\begin{theorem}\label{thm:z28}
Let $G = \mathbb{Z}_{m_1} \times \dots \times \mathbb{Z}_{m_r}$ of order $v$,
$H = \mathbb{Z}_{d_1} \times \dots \times \mathbb{Z}_{d_r}$ with $d_i \mid m_i$, $\phi$ the
reduction map and $c = v/|H|$. Let $n = k - \lambda > 0$ with $\gcd(n, |H|) = 1$,
$t \in M(H, n)$, and $s^2 = k + \lambda(v-1)$. If no $a \colon H \to \mathbb{Z}$ with
$a(th) = a(h)$ for all $h$ meets the coset conditions for $(c, k, \lambda, s)$, then $G$ has no
$(v, k, \lambda)$ SDS.
\end{theorem}

\begin{proof}
The coset sums of an SDS have autocorrelation $n[g = 0] + \lambda c$
(Lemma~\ref{lem:contraction}), so by Theorem~\ref{thm:multiplier} a translate $a$ of them is
fixed by $t$. A translate of the coset sums is the coset sums of a translate of $D$, so $a$ or
$-a$ meets the coset conditions for $(c, k, \lambda, s)$ (Definition~\ref{def:cosetconds}).
\end{proof}

\begin{remark}
Lean: \lean{not_hasSDS_of_cosetRow} (squarefree $n$) and \lean{not_hasSDS_of_cosetRow'}
(\texttt{CosetSearch.lean}, \texttt{CosetSearchWitt.lean}); both go through
\lean{not_hasSDS_of_cosetSearch}, which takes the multiplier theorem for coset sums as a
hypothesis (\lean{C10Holds}), proved by \lean{c10Holds_of_squarefree} and
\lean{c10Holds_of_ne_zero}. It settles \zrows{28} rows. Origin: the search
of~\cite[\S5]{gordon} (after Arasu, Gordon and Zhang~\cite[Algorithm~1]{arasu-gordon-zhang}),
on McFarland's theorem; its exhaustion as a kernel check and the Lean proof are ours.
\end{remark}

\subsection{Self-conjugate primes enlarge the multiplier group}

A prime $p \mid n$ restricts $t$ to powers of $p$. If $p$ is self-conjugate on $H$, the
restriction can be dropped: the coset sums are then constant modulo $p^{v_p(n)/2}$, which
already gives $p^{v_p(n)} \mid W$ for every $t$.

\begin{theorem}\label{thm:z31}
Let $D$ be an SDS $(v, k, \lambda)$ on $G$, $\phi \colon G \to H$ surjective, $c = v/|H|$, $a$
the coset sums, and $n = k - \lambda \ne 0$ with $\gcd(n, |H|) = 1$. Let $M'(H, n)$ be the set
of integers $t$ prime to $|H|$ such that for every prime $p \mid n$, either $tx = p^f x$ on $H$
for some $f \ge 0$, or $v_p(n)$ is even and $p^i x = -x$ on $H$ for some $i > 0$. Then
\begin{enumerate}
\item there is $b \in H$ with $a(th + b) = a(h + b)$ for all $h \in H$ and all
  $t \in M'(H, n)$;
\item consequently Theorem~\ref{thm:z28} holds with $M'(H, n)$ in place of $M(H, n)$.
\end{enumerate}
\end{theorem}

\begin{proof}
(1) We repeat the proof of Theorem~\ref{thm:multiplier} with $\mu = \lambda c$; only Step~1
changes, and only for a prime $p \mid n$ of the second kind, $v_p(n) = 2j$, $p$
self-conjugate on $H$. For such $p$, Turyn's lemma (Lemma~\ref{lem:turyn}, as used for
Theorem~\ref{thm:z8}) gives $a = rH + p^j Z$ with $r \in \mathbb{Z}$, $Z \in \mathbb{Z}[H]$. For any
$t$ prime to $|H|$,
\[
  W = A^{(t)}A^{(-1)} - \mu H = \kappa H + p^{2j} Z^{(t)}Z^{(-1)}, \qquad
  \kappa = r^2|H| + 2rp^j\textstyle\sum_h Z(h) - \mu .
\]
Summing coefficients, $\kappa|H| + p^{2j}(\sum_h Z(h))^2 = \sum_h W(h) = n$, so
$p^{2j} \mid \kappa |H|$, and $p \nmid |H|$ gives $p^{2j} \mid \kappa$. Hence $p^{2j} \mid W$.
The other primes are handled as in Step~1, and Steps~2--4 are unchanged.
(2) As the proof of Theorem~\ref{thm:z28}, with (1) in place of Theorem~\ref{thm:multiplier}.
\end{proof}

\begin{remark}
If $v_p(n)$ is odd, $p$ is self-conjugate on $H$ and $|H| > 1$, there is no SDS at all
($|\chi(a)|^2 = n$ for a nontrivial character $\chi$), so the parity condition loses nothing.
Lean: (1) \lean{IsSDS.exists_translate_fixed_contract_selfConj}
(\texttt{MultiplierSelfConj.lean}, with \lean{SelfConjEven}); (2)
\lean{not_hasSDS_of_cosetRowSelfConj} (\texttt{CosetSearchSelfConj.lean}) and
\lean{not_hasSDS_of_cosetRowCongr} (\texttt{CosetSearchCongr.lean}), which also restricts the
search to vectors constant modulo $p^j$, trying every residue. It settles \zrows{31} rows.
Origin: for difference sets, the self-conjugate exemption is Gordon and
Schmidt's~\cite[Thm.~3.1]{gordon-schmidt} (generalizing a theorem of Leung, Ma and Schmidt),
with the variant~\cite[Thm.~7]{miller}; self-conjugacy goes back to Turyn~\cite{turyn}. The
form for SDS coset sums on quotients prime to $n$ and its Lean proof are ours.
\end{remark}

\subsection{Field descent without ramification}

On a cyclic quotient, a large multiplier group forces the coset sums' characters to live on a
small subgroup, and Parseval bounds $n$.

\begin{theorem}\label{thm:z30}
Let $D$ be an SDS $(v, k, \lambda)$ on $G$, $n = k - \lambda \neq 0$, and
$\phi \colon G \to \mathbb{Z}_m$ surjective with $m > 1$, $v = cm$ and $\gcd(m, n) = 1$. Let
$Fd = m$, and suppose that for every $u < d$ the integer $t_u = 1 + Fu$ is prime to $m$ and
congruent modulo $m$ to a power of each prime $p \mid n$. Then
\[
  \varphi(m)\, n \le m F c^2 .
\]
\end{theorem}

\begin{proof}
Each $t_u$ lies in $M(\mathbb{Z}_m, n)$, so by Theorem~\ref{thm:multiplier} a translate $Y$ of
the coset sums is fixed by every $t_u$; it has autocorrelation $n[g=0] + \lambda c$ and
$|Y(h)| \le c$. Let $\zeta = e^{2\pi i/m}$ and $\chi_\alpha(Y) = \sum_h Y(h)\zeta^{\alpha h}$. For
$\alpha \neq 0$, $|\chi_\alpha(Y)|^2 = n$. Let $\alpha$ be a unit. Substituting $h \mapsto t_u h$,
\[
  \chi_\alpha(Y) = \sum_h Y(t_u h)\zeta^{\alpha t_u h} = \sum_h Y(h)\zeta^{\alpha h} w_h^u,
  \qquad w_h = \zeta^{\alpha F h},
\]
and averaging over $u < d$: $w_h$ is a $d$-th root of unity, equal to $1$ exactly when
$Fh = 0$ in $\mathbb{Z}_m$ ($\alpha$ is a unit), so $\sum_{u<d} w_h^u$ is $d$ or $0$. Hence
$\chi_\alpha(Y) = \chi_\alpha(Y_S)$, where $Y_S$ is $Y$ restricted to the subgroup
$S = \{h : Fh = 0\}$ of order $F$. By Parseval,
\[
  \varphi(m)\,n = \sum_{\alpha \text{ unit}} |\chi_\alpha(Y_S)|^2
  \le \sum_{\alpha \in \mathbb{Z}_m} |\chi_\alpha(Y_S)|^2 = m\sum_{h \in S} Y(h)^2 \le mFc^2. \qedhere
\]
\end{proof}

\begin{remark}
Lean: \lean{IsSDS.totient_mul_le}; rows through \lean{DescentRowCert.not_hasSDS}
(\texttt{UnramifiedDescent.lean}), where a certificate gives $m$, $F$ and the map $\phi$. It
settles \zrows{30} rows. Origin: Schmidt's field descent and the Leung--Schmidt bound, in the
form of~\cite[Results~2.11--2.12]{tan}, with McFarland's multiplier theorem and the fixed
translate of~\cite{mcfarland-rice}. Here the descent comes from the multiplier group and $F$ is
read off it; for $n$ prime to $m$ the bound is at least as strong as Schmidt's, and stronger on
N9792, $(200, 139, 18)$ in $C_2 \times C_{100}$, where it allows $F = 10$ against
$F(100, 121) = 20$ in~\cite{tan}. This form and its Lean proof are ours.
\end{remark}

\subsection{Exhaustive search}

Four Solutions settle single rows or a few rows by direct search, on $G$ or after contraction.

\begin{theorem}\label{thm:z1}
$C_3 \times C_3$ has no $(9, 8, 1)$ SDS.
\end{theorem}

\begin{proof}
Exactly one element is unmarked; by Lemma~\ref{lem:translate} we may translate it to $0$.
Lean evaluates the definition of an SDS on all $3^8$ assignments of $\{0, \pm1\}$ to the other
eight elements and finds none.
\end{proof}

\begin{remark}
Lean: \lean{Answer.N3.not_hasSDS} (case N3), by \texttt{decide}. It settles \zrows{1} row.
Origin: ours.
\end{remark}

\begin{theorem}\label{thm:z2}
$C_2 \times C_{10}$ has no $(20, 17, 8)$ SDS.
\end{theorem}

\begin{proof}
Here $s = 13$. Contract to $H = C_5$, so $c = 4$. By Lemmas~\ref{lem:contraction}
and~\ref{lem:cosetconds} the coset sums, after a sign change, would meet the coset conditions
for $(4, 17, 8, 13)$: values in $[-4, 4]$, sum $13$, square sum $41$, autocorrelation $32$.
A search of the vectors in $[-4, 4]^5$, run in Lean, finds none.
\end{proof}

\noindent By hand: for a generator $\chi$ of the characters of $C_5$, $|\chi(a)|^2 = n = 9$.
As $3$ has order $4$ modulo $5$, it is inert in $\mathbb{Z}[\zeta_5]$, so
$(\chi(a)) = (3)$ and $\chi(a) = 3u$ with $u$ a unit; the conjugates of $\chi(a)$ are the
values of the other nontrivial characters, so $u$ has absolute value $1$ in every embedding
and is a root of unity, $\chi(a) = \pm 3\zeta_5^j$. Then $a - (\pm 3[j])$ vanishes under every nontrivial
character, so $a = \pm 3[j] + rH$ with $r \in \mathbb{Z}$; the sum $13 = \pm 3 + 5r$ forces
$r = 2$ and $a(j) = 5 > c$.

\begin{remark}
Lean: \lean{Answer.N32.not_hasSDS} (case N32), by \lean{not_hasSDS_of_noVector}. It settles
\zrows{2} row. Origin: ours.
\end{remark}

\begin{theorem}\label{thm:z26}
Neither $C_3^3$ nor $C_3 \times C_9$ has a $(27, 22, 9)$ SDS.
\end{theorem}

\begin{proof}
Here $s = 16$. Contract to $H = C_3 \times C_3$, $c = 3$: the coset sums would meet the coset
conditions for $(3, 22, 9, 16)$ (values in $[-3,3]$, square sum $40$, autocorrelation $27$).
A search over $H$, run in Lean, finds no such vector.
\end{proof}

\begin{remark}
Lean: \lean{Answer.N79.not_hasSDS}, \lean{Answer.N80.not_hasSDS} (cases N79, N80), by
\lean{not_hasSDS_of_noVector}. It settles \zrows{26} rows. Origin: ours; we found
$(27, 22, 9)$ nowhere in print.
\end{remark}

When contraction alone leaves solutions, each surviving coset-sum vector can be lifted back to
$G$ coset by coset.

\begin{theorem}\label{thm:z33}
There is no SDS with parameters $(27, 17, 4)$ in $C_3^3$; $(27, 22, 3)$ in $C_3^3$ or
$C_3 \times C_9$; $(32, 28, 12)$ in $C_2^5$, $C_2^3 \times C_4$, $C_2^2 \times C_8$ or
$C_2 \times C_{16}$; $(36, 29, 20)$ in $C_3 \times C_{12}$ or $C_6 \times C_6$.
\end{theorem}

\begin{proof}
For each group take the quotient $H$ of Table~\ref{tab:lift}. \emph{Phase~1}: list every
$a \colon H \to \mathbb{Z}$ that meets the conditions of Lemmas~\ref{lem:contraction}
and~\ref{lem:cosetconds} except the marked-count bound, normalized by a translation and a
sign change so that $\sum_h a(h) = s$ and $|a(0)| = \max_h |a(h)|$; the coset sums of a
translate of $D$ or $-D$ are on the list. \emph{Phase~2}: for each listed $a$, search all
$D \colon G \to \{0, \pm1\}$ with coset sums $a$ and autocorrelation $\lambda$, assigning the
cosets one at a time in a fixed order, each by every sign pattern with the prescribed sum.
After each coset the partial autocorrelation at every shift $g$ (the sum over pairs
$(x, x+g)$ with both ends assigned) must lie within the number of pairs not yet both assigned
of its target ($k$ at $g = 0$, $\lambda$ elsewhere). A second cut uses a list of affine bijections
$\alpha$ of $G$ fixing every coset: these map SDSs with coset sums $a$ to SDSs with coset sums
$a$, so it suffices to search assignments that are lexicographically least among their images
$D \circ \alpha$; an assignment is cut as soon as an image of its assigned part is smaller. Both
phases find nothing in all nine cases.
\end{proof}

\begin{table}[ht]
\centering\small
\begin{tabular}{@{}llll@{}}
\toprule
Case & $(v, k, \lambda)$ & $G \to H$ & $c$ \\ \midrule
N70 & $(27, 17, 4)$ & $C_3^3 \to C_3^2$ & 3 \\
N76, N77 & $(27, 22, 3)$ & $C_3^3,\ C_3 \times C_9 \to C_3^2$ & 3 \\
N134, N135, N136 & $(32, 28, 12)$ & $C_2^5,\ C_2^3 \times C_4,\ C_2^2 \times C_8 \to C_2^3$ & 4 \\
N138 & $(32, 28, 12)$ & $C_2 \times C_{16} \to C_2 \times C_4$ & 4 \\
N177 & $(36, 29, 20)$ & $C_3 \times C_{12} \to C_3 \times C_4$ & 3 \\
N178 & $(36, 29, 20)$ & $C_6 \times C_6 \to C_2 \times C_6$ & 3 \\
\bottomrule
\end{tabular}
\caption{Quotients for Theorem~\ref{thm:z33}.}\label{tab:lift}
\end{table}

\begin{remark}
Lean: \lean{LiftSearch.not_hasSDS_of_lift} (\texttt{LiftSearch.lean}); rows
\lean{LiftRows.n70_not_hasSDS} and so on (\texttt{LiftRows/Part01}--\texttt{09}). It settles
\zrows{33} rows. Origin: contraction and lift as in~\cite[\S5]{gordon}; these rows and the
Lean search are ours.
\end{remark}

\subsection{How the searches are checked}\label{sec:searchcheck}

Each search above is a Boolean function written in Lean, together with a soundness theorem
proved once: if the function returns \texttt{true} on given parameters, the corresponding
hypothesis (``no $D$ fixed by $t$ \dots'', ``no $a$ \dots'') holds. A row is then closed by
evaluating the function on its parameters with \texttt{decide +kernel}: Lean's kernel reduces
the expression to \texttt{true} by its own evaluation rules, with natural-number arithmetic
done by the kernel's built-in big-number routines. No compiled code is trusted
(\texttt{native\_decide} is not used), so the only trusted component is the kernel. We describe
what a reader needs to see that the functions compute what their soundness theorems say.

\emph{Groups as numbers.} The elements of
$\mathbb{Z}_{m_1} \times \dots \times \mathbb{Z}_{m_r}$ are numbered $0, \dots, v-1$ in mixed
radix, $0$ first, and addition is computed on the numbers; a lemma proves that this numbering
is a bijection compatible with addition (\texttt{Index.lean}). Sums over $G$ thus become sums
over $0, \dots, v-1$.

\emph{Untrusted data, checked.} The orbits of $x \mapsto tx$ are found by a routine whose
correctness is not proved; instead the search checks its output. Each element $y$ gets a label
and a distance, and \lean{OrbitSearch.labelsOK} checks that every $y$ is either a root
(distance $0$, its own label) or that $ty$ has the same label and a smaller distance. By
induction on the distance, every function fixed by $t$ is constant on each label class, which
is all the soundness proof needs. Likewise the lift search's affine maps are checked by
\lean{LiftSearch.mapOK} (bijective, affine, fixing every coset), and Phase~1's list is checked
by rerunning the vector search with membership in the list as its leaf test.

\emph{The orbit search} (Theorems~\ref{thm:z15}, \ref{thm:z27}) assigns a sign or $0$ to each
class. It prunes exactly by counts: a bit mask records which pairs (number of $+$, number of
$-$) the remaining classes can still reach, and a branch survives only if the totals
$(\tfrac{k+s}{2}, \tfrac{k-s}{2})$ stay reachable. At a leaf it evaluates the autocorrelation at
one shift per nonzero class root, as $\sum_{c, c'} W_g(c, c')\,\sigma_c\sigma_{c'}$ with
precomputed weights $W_g(c, c') = \#\{y \in c : y + g \in c'\}$; the leaf must miss $\lambda$ at
some shift. (Checking only some shifts is sound: one miss suffices.)

\emph{The coset search} (Theorems~\ref{thm:z28}, \ref{thm:z31}) assigns a value in $[-c, c]$
to each class of $t$ on $H$, larger classes first. A value is tried only if the remaining
classes, of total size $R$, can still bring the sum to $s$ and the square sum to $Q$: with
remaining sum $S_r$ and square sum $Q_r$ this needs $Q_r \ge 0$ and $S_r^2 \le R\,Q_r$
(Cauchy--Schwarz) besides $|a| \le c$. A leaf must violate one of the coset conditions. In
\texttt{CosetSearchCongr.lean} the values are further restricted to one residue class modulo
$p^j$, and the search is run for every residue.

\emph{The lift search} (Theorem~\ref{thm:z33}) keeps signs in two bit masks and the partial
autocorrelation per shift as two counts; the cut of the proof above is proved sound by
bounding the change of each product term by $1$ per pair not yet assigned.

\emph{Engineering.} Integers are shifted into $\mathbb{N}$ (values $x$ stored as $x + c$, signs
coded $0, 1, 2$) so that the kernel uses its fast natural-number arithmetic; lemmas transfer the
results back. Large searches are split: one declaration checks that the tree has at most $N$
nodes at depth $d$, others check the subtrees below slices of these nodes, and a theorem joins
them (\texttt{OrbitSearchSplit.lean}, \lean{cosetSearch_of_slices},
\lean{LiftSearch.liftRow_of_slices}). The arithmetic hypotheses of each rule (primality,
squarefreeness, coprimality, $s^2 = k + \lambda(v-1)$, $t \equiv p^f$) are decided in the same
way, by trial division.

\emph{Controls.} The Lean files run the same functions on parameters where an SDS exists, for
instance $(11, 6, 1)$ in $C_{11}$ and $(19, 10, 3)$ in $C_{19}$, and check that the search
returns \texttt{false} there. Independent searches by the refuter, written separately from the
Lean code, agree on the rows they recomputed (Reviews R26, R27, R30, R31, R32).
